\documentclass[../../main.tex]{subfiles}

\begin{document}

\section{关于 $L^{p}(1\leqslant p<\infty)$ 的对偶的 Riesz 表示定理}\label{section:ch08-Riesz表示}

\begin{definition}[线性泛函]\label{definition:ch08-线性泛函}
线性空间 $X$ 上的线性泛函是 $X$ 上的实值函数 $T$, 它使得对 $X$ 中的函数 $g$ 和 $h$ 以及实数 $\alpha$ 和 $\beta$,
\begin{align*}
T(\alpha\cdot g+\beta\cdot h)=\alpha\cdot T(g)+\beta\cdot T(h).
\end{align*}
\end{definition}

容易验证逐点定义的线性泛函的线性组合也是线性的. 因此线性空间上的线性泛函组成的族自身是一个线性空间.

\begin{example}\label{example:ch08-积分泛函}
令 $E$ 为可测集, $1\leqslant p<\infty$, $q$ 为 $p$ 的共轭, 而 $g$ 属于 $L^{q}(E)$. 定义 $L^{p}(E)$ 上的泛函 $T$ 为
\begin{align}
\text{对所有 }f\in L^{p}(E),\quad T(f)=\int_{E}g\cdot f.\label{align:ch08-积分泛函}
\end{align}
H\"older 不等式告诉我们对于 $f\in L^{p}(E)$, 乘积 $g\cdot f$ 在 $E$ 上可积, 因此泛函 $T$ 是正确定义的. 根据积分的线性, $T$ 是线性的. 观察到 H\"older 不等式说的是对所有 $f\in L^{p}(E)$,
\begin{align}
|T(f)|\leqslant\|g\|_{q}\cdot\|f\|_{p}.\label{align:ch08-积分泛函估计}
\end{align}
\end{example}

\begin{example}\label{example:ch08-Riemann-Stieltjes}
令 $[a,b]$ 为闭有界区间而函数 $g$ 在 $[a,b]$ 上是有界变差的. 定义 $C[a,b]$ 上的泛函 $T$ 为
\begin{align}
\text{对所有 }f\in C[a,b],\quad T(f)=\int_{a}^{b}f(x)\,dg(x),\label{align:ch08-Stieltjes}
\end{align}
其中积分是 Riemann-Stieltjes 意义下的积分. 泛函 $T$ 是正确定义的且是线性的\footnote{关于 Riemann-Stieltjes 积分见 Richard Wheedon 和 Antoni Zygmund 的书 \emph{Measure and Integral}[WZ77] 的第 2 章.}. 此外, 从积分的定义立即得出, 对所有 $f\in C[a,b]$,
\begin{align}
|T(f)|\leqslant TV(g)\cdot\|f\|_{\max},\label{align:ch08-Stieltjes估计}
\end{align}
其中 $TV(g)$ 是 $g$ 在 $[a,b]$ 上的全变差.
\end{example}

\begin{definition}[有界线性泛函与范数]\label{definition:ch08-有界泛函}
对于赋范线性空间 $X$, $X$ 上的线性泛函 $T$ 称为是有界的, 若存在 $M\geqslant 0$, 使得
\begin{align}
\text{对所有 }f\in X,\quad|T(f)|\leqslant M\cdot\|f\|.\label{align:ch08-有界泛函}
\end{align}
所有这样的 $M$ 的下确界称为 $T$ 的范数, 记为 $\|T\|_{*}$.
\end{definition}

不等式\eqref{align:ch08-积分泛函估计}告诉我们第一个例子中的线性泛函是有界的, 而不等式\eqref{align:ch08-Stieltjes估计}证明第二个例子也是有界的.

令 $T$ 为赋范线性空间 $X$ 上的有界线性泛函. 容易看出\eqref{align:ch08-有界泛函}对 $M=\|T\|_{*}$ 成立. 因此, 根据 $T$ 的线性, 对所有 $f$, $h\in X$,
\begin{align}
|T(f)-T(h)|\leqslant\|T\|_{*}\cdot\|f-h\|.\label{align:ch08-泛函连续}
\end{align}
由此我们推出有界线性泛函 $T$ 的以下连续性质:
\begin{align}
\text{若在 }X\text{ 中 }\{f_{n}\}\to f,\quad\text{则}\{T(f_{n})\}\to T(f).\label{align:ch08-泛函连续性}
\end{align}
我们把以下证明留作练习: 证明
\begin{align}
\|T\|_{*}=\sup\{T(f)\mid f\in X,\quad\|f\|\leqslant 1\}.\label{align:ch08-范数刻画}
\end{align}
且用 $\|\cdot\|_{*}$ 的这个刻画去证明下面的命题.

\begin{proposition}\label{proposition:ch08-对偶空间}
令 $X$ 为赋范线性空间. 则 $X$ 上的有界线性泛函族是一个线性空间, 在其上 $\|\cdot\|_{*}$ 是一个范数. 该赋范线性空间称为 $X$ 的对偶空间且记为 $X^{*}$.
\end{proposition}

\begin{proposition}\label{proposition:ch08-积分泛函范数}
令 $E$ 为可测集, $1\leqslant p<\infty$, $q$ 为 $p$ 的共轭, 而 $g$ 属于 $L^{q}(E)$. 定义 $L^{p}(E)$ 上的泛函 $T$ 为
\begin{align*}
\text{对所有 }f\in L^{p}(E),\quad T(f)=\int_{E}g\cdot f.
\end{align*}
则 $T$ 是 $L^{p}(E)$ 上的有界线性泛函且 $\|T\|_{*}=\|g\|_{q}$.
\end{proposition}

\begin{proof}
我们从\eqref{align:ch08-积分泛函估计}推出 $T$ 是 $L^{p}(E)$ 上的一个有界线性泛函且 $\|T\|_{*}\leqslant\|g\|_{q}$. 另一方面, 对于 $p>1$, 根据前一章的定理 1 ($p$ 与 $q$ 互换), $g$ 的共轭函数 $g^{*}=\|g\|_{q}^{1-q}\mathrm{sgn}(g)\cdot|g|^{q-1}$ 属于 $L^{p}(E)$,
\begin{align*}
T(g^{*})=\|g\|_{q}\text{ 且 }\|g^{*}\|_{p}=1.
\end{align*}
从\eqref{align:ch08-范数刻画}得出 $\|T\|_{*}=\|g\|_{q}$. 对于 $p=1$, 我们用反证法. 若 $\|g\|_{\infty}>\|T\|_{*}$, 则存在具有有限测度的集合 $A$ 使得在其上 $|g|>\|T\|_{*}$. 定义 $f=[1/m(A)][\mathrm{sgn}g]\chi_{A}$. 则 $\|f\|_{1}=1$, 然而 $T(f)>\|T\|_{*}$, 这是一个矛盾. 从而得出 $p=1$ 时, $\|T\|_{*}=\|g\|_{\infty}$.
\end{proof}

现在我们的目标是证明对于 $1\leqslant p<\infty$, $L^{p}(E)$ 上的每个有界线性泛函由 $L^{q}(E)$ 中的函数的积分给出, 其中 $q$ 是 $p$ 的共轭.

\begin{proposition}\label{proposition:ch08-稠密集相等}
令 $T$ 和 $S$ 为赋范线性空间 $X$ 上的有界线性泛函. 若在 $X$ 的一个稠密子集 $X_{0}$ 上 $T=S$, 则 $T=S$.
\end{proposition}

\begin{proof}
令 $g$ 属于 $X$. 由于 $X_{0}$ 在 $X$ 中稠密, 存在 $X_{0}$ 中的序列 $\{g_{n}\}$ 在 $X$ 中收敛于 $g$. 我们从\eqref{align:ch08-泛函连续性}推出 $\{S(g_{n})\}\to S(g)$ 以及 $\{T(g_{n})\}\to T(g)$. 但对所有 $n$, $S(g_{n})=T(g_{n})$, 因此 $S(g)=T(g)$.
\end{proof}

\begin{lemma}\label{lemma:ch08-积分不等式推出Lq}
令 $E$ 为可测集而 $1\leqslant p<\infty$. 假定函数 $g$ 在 $E$ 上可积且存在 $M\geqslant 0$ 使得
\begin{align}
\text{对 }L^{p}(E)\text{ 中的每个简单函数 }f,\quad\left|\int_{E}g\cdot f\right|\leqslant M\|f\|_{p}.\label{align:ch08-简单函数不等式}
\end{align}
则 $g$ 属于 $L^{q}(E)$, 其中 $q$ 是 $p$ 的共轭. 此外, $\|g\|_{q}\leqslant M$.
\end{lemma}

\begin{proof}
由于 $g$ 在 $E$ 上是可积的, 它在 $E$ 上 a.e. 有限. 通过可能地从 $E$ 中去掉测度为零的集合, 我们假设 $g$ 在整个 $E$ 上有限. 我们首先考虑 $p>1$ 的情形. 由于 $|g|$ 是非负可测函数, 根据简单逼近定理, 存在在 $E$ 上逐点收敛于 $|g|$ 且在 $E$ 上对所有 $n$ 有 $0\leqslant\varphi_{n}\leqslant|g|$ 的简单函数序列 $\{\varphi_{n}\}$. 由于 $\{\varphi_{n}^{q}\}$ 是在 $E$ 上逐点收敛于 $|g|^{q}$ 的非负可测函数序列, 根据 Fatou 引理, 要证明 $|g|^{q}$ 在 $E$ 上可积且 $\|g\|_{q}\leqslant M$, 仅须证明对所有 $n$,
\begin{align}
\int_{E}\varphi_{n}^{q}\leqslant M^{q}.\label{align:ch08-简单函数q积分}
\end{align}
固定自然数 $n$. 为证明\eqref{align:ch08-简单函数q积分}我们估计 $\varphi_{n}^{q}$ 在 $E$ 上的泛函值如下: 在 $E$ 上,
\begin{align}
\varphi_{n}^{q}=\varphi_{n}\cdot\varphi_{n}^{q-1}\leqslant|g|\cdot\varphi_{n}^{q-1}=g\cdot\mathrm{sgn}(g)\cdot\varphi_{n}^{q-1}.\label{align:ch08-简单函数q估计}
\end{align}
我们定义简单函数 $f_{n}$ 为, 在 $E$ 上,
\begin{align*}
f_{n}=\mathrm{sgn}(g)\cdot\varphi_{n}^{q-1}.
\end{align*}
函数 $\varphi_{n}$ 在 $E$ 上是可积的, 这是由于它在 $E$ 上被可积函数 $g$ 控制. 由于 $\varphi_{n}$ 是简单的, 它有有限支撑, 因此 $f_{n}$ 属于 $L^{p}(E)$. 我们从\eqref{align:ch08-简单函数q估计}和\eqref{align:ch08-简单函数不等式}推出,
\begin{align}
\int_{E}\varphi_{n}^{q}\leqslant\int_{E}g\cdot f_{n}\leqslant M\|f_{n}\|_{p}.\label{align:ch08-简单函数q估计2}
\end{align}
由于 $q$ 是 $p$ 的共轭, $p(q-1)=q$, 因此,
\begin{align*}
\int_{E}|f_{n}|^{p}=\int_{E}\varphi_{n}^{p(q-1)}=\int_{E}\varphi_{n}^{q}.
\end{align*}
我们重写\eqref{align:ch08-简单函数q估计2}为
\begin{align*}
\int_{E}\varphi_{n}^{q}\leqslant M\cdot\left[\int_{E}\varphi_{n}^{q}\right]^{1/q}.
\end{align*}
由于 $\varphi_{n}^{q}$ 在 $E$ 上是可积的, 我们可以重新整理该积分不等式为
\begin{align*}
\left[\int_{E}\varphi_{n}^{q}\right]^{1-1/p}\leqslant M.
\end{align*}
由于 $1-1/p=1/q$, 它是\eqref{align:ch08-简单函数q积分}的重述.

剩下来考虑 $p=1$ 的情形. 我们必须证明 $M$ 是 $g$ 的一个本性上界. 我们用反证法. 若 $M$ 不是本性上界, 则根据测度的连续性, 存在某个 $\varepsilon>0$, 使得集合 $E_{\varepsilon}=\{x\in E\mid|g(x)|>M+\varepsilon\}$ 的测度不为零. 若令 $f$ 为 $E_{\varepsilon}$ 的具有有限正测度的可测子集的特征函数, 我们得到与\eqref{align:ch08-简单函数不等式}矛盾的结论.
\end{proof}

\begin{theorem}\label{theorem:ch08-Riesz区间}
令 $[a,b]$ 为闭有界区间而 $1\leqslant p<\infty$. 假定 $T$ 是 $L^{p}[a,b]$ 上的有界线性泛函. 则存在 $L^{q}[a,b]$ 中的函数 $g$, 其中 $q$ 是 $p$ 的共轭, 使得对所有 $L^{p}[a,b]$ 中的 $f$,
\begin{align*}
T(f)=\int_{a}^{b}g\cdot f.
\end{align*}
\end{theorem}

\begin{proof}
我们考虑 $p>1$ 的情形. $p=1$ 的情形的证明是类似的. 对 $[a,b]$ 中的 $x$, 定义
\begin{align*}
\Phi(x)=T(\chi_{[a,x]}).
\end{align*}
我们宣称这个实值函数 $\Phi$ 在 $[a,b]$ 上是绝对连续的. 事实上, 根据 $T$ 的线性, 对每个 $[c,d]\subseteq[a,b]$, 由于 $\chi_{[c,d]}=\chi_{[a,d]}-\chi_{[a,c]}$,
\begin{align*}
\Phi(d)-\Phi(c)=T(\chi_{[a,d]})-T(\chi_{[a,c]})=T(\chi_{[c,d]}).
\end{align*}
因此, 若 $\{(a_{k},b_{k})\}_{k=1}^{n}$ 是 $(a,b)$ 中的区间的有限不交族, 根据 $T$ 的线性,
\begin{align}
\sum_{k=1}^{n}|\Phi(b_{k})-\Phi(a_{k})|=\sum_{k=1}^{n}\varepsilon_{k}\cdot T(\chi_{[a_{k},b_{k}]})=T\left(\sum_{k=1}^{n}\varepsilon_{k}\cdot\chi_{[a_{k},b_{k}]}\right),\label{align:ch08-绝对连续估计}
\end{align}
其中每个 $\varepsilon_{k}=\mathrm{sgn}[\Phi(b_{k})-\Phi(a_{k})]$. 此外, 对于简单函数 $f=\sum_{k=1}^{n}\varepsilon_{k}\cdot\chi_{[a_{k},b_{k}]}$,
\begin{align*}
|T(f)|\leqslant\|T\|_{*}\cdot\|f\|_{p}\text{ 且 }\|f\|_{p}=\left[\sum_{k=1}^{n}(b_{k}-a_{k})\right]^{1/p}.
\end{align*}
于是,
\begin{align*}
\sum_{k=1}^{n}|\Phi(b_{k})-\Phi(a_{k})|\leqslant\|T\|_{*}\cdot\left[\sum_{k=1}^{n}(b_{k}-a_{k})\right]^{1/p}.
\end{align*}
因此, $\delta=(\varepsilon/\|T\|_{*})^{p}$ 对应 $\Phi$ 在 $[a,b]$ 上绝对连续准则的任何 $\varepsilon>0$ 挑战.

根据第 6 章的定理 10, 函数 $g=\Phi'$ 在 $[a,b]$ 上是可积的且对所有 $x\in[a,b]$,
\begin{align*}
\Phi(x)=\int_{a}^{x}g.
\end{align*}
因此, 对每个 $[c,d]\subseteq(a,b)$,
\begin{align*}
T(\chi_{[c,d]})=\Phi(d)-\Phi(c)=\int_{a}^{b}g\cdot\chi_{[c,d]}.
\end{align*}
由于泛函 $T$\footnote{泛函 $T$ 必须遵循 $L^{p}[a,b]$ 中的函数在 $[a,b]$ 上 a.e. 相等的等价关系, 特别地, 对于 $a\leqslant c\leqslant d\leqslant b$, $T(\chi_{[c,d]})=T(\chi_{(c,d)})=T(\chi_{[c,d)})$.}和泛函 $f\mapsto\int_{a}^{b}g\cdot f$ 在阶梯函数的线性空间上是线性的, 这就得到
\begin{align*}
\text{对 }[a,b]\text{ 上的所有阶梯函数 }f,\quad T(f)=\int_{a}^{b}g\cdot f.
\end{align*}
根据前一章的命题 10 及其证明, 若 $f$ 是 $[a,b]$ 上的简单函数, 则存在在 $L^{p}[a,b]$ 中收敛到 $f$ 的阶梯函数序列 $\{\varphi_{n}\}$, 它在 $[a,b]$ 上也是一致逐点有界的. 由于线性泛函 $T$ 在 $L^{p}[a,b]$ 上是有界的, 从连续性\eqref{align:ch08-泛函连续性}得到
\begin{align*}
\lim_{n\to\infty}T(\varphi_{n})=T(f).
\end{align*}
另一方面, 根据 Lebesgue 控制收敛定理,
\begin{align*}
\lim_{n\to\infty}\int_{a}^{b}g\cdot\varphi_{n}=\int_{a}^{b}g\cdot f.
\end{align*}
因此,
\begin{align*}
\text{对 }[a,b]\text{ 上的所有简单函数 }f,\quad T(f)=\int_{a}^{b}g\cdot f.
\end{align*}
由于 $T$ 是有界的, 对 $[a,b]$ 上的所有简单函数 $f$,
\begin{align*}
\left|\int_{a}^{b}g\cdot f\right|=|T(f)|\leqslant\|T\|_{*}\cdot\|f\|_{p}.
\end{align*}
根据引理 4, $g$ 属于 $L^{q}[a,b]$. 从命题 2 得出线性泛函 $f\mapsto\int_{a}^{b}g\cdot f$ 在 $L^{p}[a,b]$ 上是有界的. 该泛函与有界泛函 $T$ 在全体简单函数上一致, 根据前一章的命题 9, 全体简单函数是 $L^{p}[a,b]$ 的一个稠密集空间. 我们从命题 3 推出这两个泛函在整个 $L^{p}[a,b]$ 上一致.
\end{proof}

\begin{theorem}[Riesz 表示定理]\label{theorem:ch08-Riesz表示}
令 $E$ 为可测集, $1\leqslant p<\infty$, 而 $q$ 为 $p$ 的共轭. 对每个 $g\in L^{q}(E)$, 定义 $L^{p}(E)$ 上的有界线性泛函 $\mathcal{R}_{g}$ 为对 $L^{p}(E)$ 中的所有 $f$,
\begin{align}
\mathcal{R}_{g}(f)=\int_{E}g\cdot f.\label{align:ch08-Riesz映射}
\end{align}
则对每个 $L^{p}(E)$ 上的有界线性泛函 $T$, 存在唯一的函数 $g\in L^{q}(E)$, 使得
\begin{align}
\mathcal{R}_{g}=T\text{ 且 }\|T\|_{*}=\|g\|_{q}.\label{align:ch08-Riesz结论}
\end{align}
\end{theorem}

\begin{proof}
命题 2 告诉我们对每个 $g\in L^{q}(E)$, $\mathcal{R}_{g}$ 是 $L^{p}(E)$ 上使得 $\|\mathcal{R}_{g}\|_{*}=\|g\|_{q}$ 的有界线性泛函. 根据积分的线性, 对每个 $g_{1}$, $g_{2}\in L^{q}(E)$,
\begin{align*}
\mathcal{R}_{g_{1}}-\mathcal{R}_{g_{2}}=\mathcal{R}_{g_{1}-g_{2}}.
\end{align*}
因此, 若 $\mathcal{R}_{g_{1}}=\mathcal{R}_{g_{2}}$, 则 $\mathcal{R}_{g_{1}-g_{2}}=0$. 于是 $\|g_{1}-g_{2}\|_{q}=0$, 从而 $g_{1}=g_{2}$. 因此, 对 $L^{p}(E)$ 上的有界线性泛函 $T$, 至多存在一个函数 $g\in L^{q}(E)$ 使得 $\mathcal{R}_{g}=T$. 剩下来要证明对每个 $L^{p}(E)$ 上的有界线性泛函 $T$, 存在函数 $g\in L^{q}(E)$ 使得 $T=\mathcal{R}_{g}$. 前一个定理告诉我们, 对 $E$ 是闭有界区间是如此. 我们现在对 $E=\mathbf{R}$ 验证这个结论且接着对一般的可测集 $E$ 验证.

令 $T$ 为 $L^{p}(\mathbf{R})$ 上的有界线性泛函. 固定一个自然数 $n$. 定义 $L^{p}[-n,n]$ 上的线性泛函 $T_{n}$ 为对所有 $f\in L^{p}[-n,n]$,
\begin{align*}
T_{n}(f)=T(\widehat{f}),
\end{align*}
其中 $\widehat{f}$ 是 $f$ 在整个 $\mathbf{R}$ 的延拓, 它在 $[-n,n]$ 外等于零. 则由于 $\|f\|_{p}=\|\widehat{f}\|_{p}$,
\begin{align*}
\text{对所有 }f\in L^{p}[-n,n],\quad|T_{n}(f)|\leqslant\|T\|_{*}\|f\|_{p}.
\end{align*}
因此 $\|T_{n}\|_{*}\leqslant\|T\|_{*}$. 前一个定理和命题 2 告诉我们存在一个函数 $g_{n}\in L^{q}[-n,n]$ 使得
\begin{align}
\text{对所有 }f\in L^{p}[-n,n],\quad T_{n}(f)=\int_{-n}^{n}g_{n}\cdot f\text{ 且 }\|g_{n}\|_{q}=\|T_{n}\|_{*}\leqslant\|T\|_{*}.\label{align:ch08-Riesz逐步}
\end{align}
根据证明开始时关于唯一性的注记, $g_{n+1}$ 在 $[-n,n]$ 上的限制与 $g_{n}$ 在 $[-n,n]$ 上的限制 a.e. 一致. 定义 $g$ 为 $\mathbf{R}$ 上的可测函数, 使得对每个 $n$ 与 $g_{n}$ 在 $[-n,n]$ 上 a.e. 一致. 我们从 $T_{n}$ 和 $g_{n}$ 的定义以及\eqref{align:ch08-Riesz逐步}左边的等式推出, 对所有在有界集外消失的函数 $f\in L^{p}(\mathbf{R})$,
\begin{align*}
T(f)=\int_{\mathbf{R}}g\cdot f.
\end{align*}
根据\eqref{align:ch08-Riesz逐步}右边的不等式, 对所有 $n$,
\begin{align*}
\int_{-n}^{n}|g|^{q}\leqslant(\|T\|_{*})^{q}.
\end{align*}
因此, 根据 Fatou 引理, $g$ 属于 $L^{q}(\mathbf{R})$. 由于有界线性泛函 $\mathcal{R}_{g}$ 和 $T$ 在 $L^{p}(\mathbf{R})$ 的稠密集空间 (该子空间由在有界集外等于零的 $L^{p}(\mathbf{R})$ 函数组成) 上一致, 从命题 3 得出 $\mathcal{R}_{g}$ 在整个 $L^{p}(\mathbf{R})$ 上与 $T$ 一致.

最后, 考虑一般的可测集 $E$, 而 $T$ 是 $L^{p}(E)$ 上的有界线性泛函. 定义 $L^{p}(\mathbf{R})$ 上的线性泛函 $\widehat{T}$ 为 $\widehat{T}(f)=T(f|_{E})$. 则 $\widehat{T}$ 是 $L^{p}(\mathbf{R})$ 上的有界线性泛函. 我们刚才已证明存在函数 $\widehat{g}\in L^{q}(\mathbf{R})$ 使得 $\widehat{T}$ 表示 $\mathbf{R}$ 上对应 $\widehat{g}$ 的积分. 定义 $g$ 为 $\widehat{g}$ 在 $E$ 上的限制. 则 $T=\mathcal{R}_{g}$.
\end{proof}

\begin{remark}\label{remark:ch08-Riesz区间注记}
在本节的第二个例子中, 我们展示了有界变差函数的积分对应 $C[a,b]$ 上的有界线性泛函. 我们将在第 21 章证明的 Riesz 的一个定理告诉我们, $C[a,b]$ 上的所有有界线性泛函都具有这种形式. 在 21.5 节, 我们刻画了 $C(K)$ 上的有界线性泛函, $C(K)$ 是赋予最大值范数的紧拓扑空间 $K$ 上的连续实值函数的线性空间.
\end{remark}

\begin{remark}\label{remark:ch08-Linfty对偶注记}
令 $[a,b]$ 为非退化的闭有界区间. 我们从积分的线性与 H\"older 不等式推出, 若 $f$ 属于 $L^{1}[a,b]$, 则泛函 $g\mapsto\int_{a}^{b}f\cdot g$ 是 $L^{\infty}[a,b]$ 上的有界线性泛函. 然而, 存在不是这种形式的 $L^{\infty}[a,b]$ 上的有界线性泛函. 在 19.3 节, 我们证明 Kantorovich 的一个定理, 它刻画了 $L^{\infty}$ 的对偶空间.
\end{remark}

\begin{exercise}\label{exercise:ch08-01}
验证\eqref{align:ch08-范数刻画}.
\end{exercise}

\begin{exercise}\label{exercise:ch08-02}
证明命题 1.
\end{exercise}

\begin{exercise}\label{exercise:ch08-03}
令 $T$ 为赋范线性空间 $X$ 上的线性泛函. 证明 $T$ 是有界的当且仅当连续性质\eqref{align:ch08-泛函连续性}成立.
\end{exercise}

\begin{exercise}\label{exercise:ch08-04}
赋范线性空间 $X$ 上的泛函 $T$ 称为是 Lipschitz 的, 若存在 $c\geqslant 0$ 使得
\begin{align*}
\text{对所有 }g,h\in X,\quad|T(g)-T(h)|\leqslant c\|g-h\|.
\end{align*}
这样的 $c$ 的下确界称为 $T$ 的 Lipschitz 常数. 证明一个线性泛函是有界的当且仅当它是 Lipschitz 的, 在这种情形它的 Lipschitz 常数是 $\|T\|_{*}$.
\end{exercise}

\begin{exercise}\label{exercise:ch08-05}
令 $E$ 为可测集而 $1\leqslant p<\infty$. 证明 $L^{p}(E)$ 中在有界集外消失的函数在 $L^{p}(E)$ 中稠密. 证明这对 $L^{\infty}(\mathbf{R})$ 不成立.
\end{exercise}

\begin{exercise}\label{exercise:ch08-06}
在 $p=1$ 的情形下证明 Riesz 表示定理. 首先证明, 沿用定理证明的记号, 函数 $\Phi$ 是 Lipschitz 的, 因此是绝对连续的. 接着遵循 $p>1$ 的证明.
\end{exercise}

\begin{exercise}\label{exercise:ch08-07}
对 $\ell^{p}(1\leqslant p<\infty)$ 空间上的有界线性泛函叙述和证明 Riesz 表示定理.
\end{exercise}

\begin{exercise}\label{exercise:ch08-08}
令 $c$ 为收敛到实数的实序列的线性空间, 而 $c_{0}$ 是收敛到 $0$ 的序列组成的 $c$ 的子空间. 赋予这些线性空间以 $\ell^{\infty}$ 范数. 确定 $c$ 和 $c_{0}$ 的对偶空间.
\end{exercise}

\begin{exercise}\label{exercise:ch08-09}
令 $[a,b]$ 为闭有界区间而 $C[a,b]$ 赋予最大值范数. 令 $x_{0}$ 属于 $[a,b]$. 定义 $C[a,b]$ 上的线性泛函 $T$ 为 $T(f)=f(x_{0})$. 证明 $T$ 是有界的且由对应有界变差函数的 Riemann-Stieltjes 积分给出.
\end{exercise}

\begin{exercise}\label{exercise:ch08-10}
令 $f$ 属于 $C[a,b]$. 证明存在 $[a,b]$ 上的有界变差函数 $g$ 使得
\begin{align*}
\int_{a}^{b}f\,dg=\|f\|_{\max}\text{ 且 }TV(g)=1.
\end{align*}
\end{exercise}

\begin{exercise}\label{exercise:ch08-11}
令 $[a,b]$ 为闭有界区间而 $C[a,b]$ 赋予最大值范数. $T$ 为 $C[a,b]$ 上的有界线性泛函. 对于 $x\in[a,b]$, 令 $g_{x}$ 为 $C[a,b]$ 的成员, 它在 $[a,x]$ 和 $[x,b]$ 上是线性的, 其中 $g_{x}(a)=0$, $g_{x}(x)=x-a$, 而 $g_{x}(b)=x-b$. 对 $x\in[a,b]$ 定义 $\Phi(x)=T(g_{x})$. 证明 $\Phi$ 在 $[a,b]$ 上是 Lipschitz 的.
\end{exercise}

\end{document}